\documentclass[a4paper,11pt]{article}
\usepackage[margin=1in]{geometry}
\usepackage{amsmath,amsthm,amssymb}

\newtheorem{theorem}{Theorem}
\newtheorem{corollary}[theorem]{Corollary}

\title{Transferring Lipschitz Tail Bounds by a Finite Coupling}
\author{}
\date{}

\begin{document}
\maketitle

\begin{abstract}
A coupling of two probability weights on the same finite metric space
describes how mass is paired between them. If all but an exceptional mass
is paired at distance at most a given radius, every real-valued
1-Lipschitz observable inherits a tail bound: the radius is enlarged by
that distance and the tail probability increases by at most the
exceptional mass. We prove this estimate directly, derive a radius--error
tradeoff from the mean coupled distance, and use a two-point example to
show why a small mean distance does not imply exact containment of the
observable values in a small interval.
\end{abstract}

\section{A tail estimate on a common metric space}

Suppose that two probability weights assign different masses to the same
finite metric space. A real-valued function may be concentrated near a
chosen center under the first weight. To transfer that information to the
second weight, we pair the masses and control how far the pairs lie
apart. A Lipschitz function then changes little on the close pairs; the
remaining pairs account for an additive error in probability.

Let $(X,d)$ be a nonempty finite metric space. A probability weight on
$X$ is a function $\mu:X\to[0,1]$ with
$\sum_{x\in X}\mu(x)=1$. For a subset $A\subseteq X$, write
$\mu(A)=\sum_{x\in A}\mu(x)$. Let $\mu$ and $\nu$ be two such weights.
A \emph{coupling} of $\mu$ and $\nu$ is an array
$\pi:X\times X\to[0,\infty)$ satisfying
\begin{equation}\label{eq:marginals}
 \sum_{y\in X}\pi(x,y)=\mu(x)\quad(x\in X),
 \qquad
 \sum_{x\in X}\pi(x,y)=\nu(y)\quad(y\in X).
\end{equation}
Thus $\pi$ has total mass one. For $E\subseteq X\times X$, put
$\pi(E)=\sum_{(x,y)\in E}\pi(x,y)$.

A function $f:X\to\mathbb R$ is \emph{1-Lipschitz} if
$|f(x)-f(y)|\le d(x,y)$ for all $x,y\in X$. The following theorem
transfers a tail bound for every such function using the same coupling.

\begin{theorem}\label{thm:transfer}
Let $\pi$ be a coupling of $\mu$ and $\nu$. Suppose that $r\ge0$ and
$0\le\eta\le1$ satisfy
\[
 \pi\bigl(\{(x,y)\in X\times X:d(x,y)>r\}\bigr)\le\eta.
\]
Then, for every 1-Lipschitz function $f:X\to\mathbb R$, every center
$a\in\mathbb R$, and every $t\ge0$,
\begin{equation}\label{eq:transfer}
 \nu\bigl(\{y\in X:|f(y)-a|>t+r\}\bigr)
 \le
 \min\left\{1,
 \mu\bigl(\{x\in X:|f(x)-a|>t\}\bigr)+\eta\right\}.
\end{equation}
\end{theorem}

The two costs in \eqref{eq:transfer} have different meanings. The radius
increment $r$ covers the change of $f$ along close pairs. The exceptional
mass $\eta$ covers the pairs on which that change has no such bound.
The center $a$ is an arbitrary real number; the theorem assumes no
centering property under either weight.

\section{Proof by an inclusion of events}

\begin{proof}[Proof of Theorem~\ref{thm:transfer}]
Fix $f$, $a$, and $t$ as in the theorem. If $x,y\in X$ satisfy
$|f(x)-a|\le t$ and $d(x,y)\le r$, then
\[
 |f(y)-a|
 \le |f(y)-f(x)|+|f(x)-a|
 \le d(x,y)+t
 \le r+t.
\]
Consequently, a pair whose second coordinate lies in the target tail
must either have its first coordinate in the original tail or lie farther
apart than $r$. In set notation,
\begin{align*}
 &\{(x,y)\in X\times X:|f(y)-a|>t+r\}\\
 &\quad\subseteq
 \{(x,y)\in X\times X:|f(x)-a|>t\}
 \ \cup\ \{(x,y)\in X\times X:d(x,y)>r\}.
\end{align*}
Summing the nonnegative weights $\pi(x,y)$ over this inclusion gives
\begin{align*}
 \nu\bigl(\{y:|f(y)-a|>t+r\}\bigr)
 &=\sum_{\substack{y\in X\\|f(y)-a|>t+r}}
       \sum_{x\in X}\pi(x,y)\\
 &\le
 \sum_{\substack{x\in X\\|f(x)-a|>t}}
       \sum_{y\in X}\pi(x,y)
 +\sum_{\substack{x,y\in X\\d(x,y)>r}}\pi(x,y)\\
 &\le \mu\bigl(\{x:|f(x)-a|>t\}\bigr)+\eta.
\end{align*}
The equality uses the column sums in \eqref{eq:marginals}, and the last
inequality uses the row sums and the assumed exceptional-mass bound.
The left side is also at most one because $\nu$ is a probability weight.
Combining these two bounds proves \eqref{eq:transfer}.
\end{proof}

\section{Using the mean coupled distance}

The exceptional-mass condition can be obtained from the mean distance of
the pairs. For a coupling $\pi$, define
\[
 D=\sum_{x,y\in X}\pi(x,y)d(x,y).
\]
The number $D$ is finite and nonnegative. It describes an average; it
does not bound the distance of every pair with positive coupling mass.

\begin{corollary}\label{cor:mean}
Let $\pi$ be a coupling of $\mu$ and $\nu$ with mean distance $D$.
For every $r>0$, every 1-Lipschitz function $f:X\to\mathbb R$, every
$a\in\mathbb R$, and every $t\ge0$,
\begin{equation}\label{eq:mean}
 \nu\bigl(\{y:|f(y)-a|>t+r\}\bigr)
 \le
 \min\left\{1,
 \mu\bigl(\{x:|f(x)-a|>t\}\bigr)
 +\min\{1,D/r\}\right\}.
\end{equation}
If $D=0$, then $\mu=\nu$.
\end{corollary}

\begin{proof}
For $r>0$, every term in the sum over pairs at distance greater than $r$
has $r\pi(x,y)\le d(x,y)\pi(x,y)$. Hence
\[
 r\,\pi\bigl(\{(x,y):d(x,y)>r\}\bigr)
 \le \sum_{\substack{x,y\in X\\d(x,y)>r}}\pi(x,y)d(x,y)
 \le D.
\]
Division by $r$ and the total mass one give an exceptional-mass bound
$\min\{1,D/r\}$. Applying Theorem~\ref{thm:transfer} with this value of
$\eta$ proves \eqref{eq:mean}.

If $D=0$, each nonnegative summand $\pi(x,y)d(x,y)$ is zero.
For $x\ne y$, the metric property gives $d(x,y)>0$, so $\pi(x,y)=0$.
The coupling is therefore supported on the diagonal, and its marginals
satisfy $\mu(x)=\pi(x,x)=\nu(x)$ for every $x\in X$.
\end{proof}

For an explicit error budget, suppose that $D>0$ and choose
$q\in(0,1]$. Setting $r=D/q>0$ in \eqref{eq:mean} yields
\begin{equation}\label{eq:budget}
 \nu\bigl(\{y:|f(y)-a|>t+D/q\}\bigr)
 \le \min\left\{1,
 \mu\bigl(\{x:|f(x)-a|>t\}\bigr)+q\right\}.
\end{equation}
Thus a smaller allowed exceptional mass requires a larger radius
increment. If the original tail has mass at most $p\in[0,1]$, the new
tail has mass at most $\min\{1,p+q\}$. When $D=0$, the equality of
weights gives equality of their tails at every radius, and no division
or radius increment is needed.

\section{A distant atom and the scope of the estimate}

Let $X=\{0,10\}$ with the usual distance, and choose $0<\varepsilon<1$.
Put $\mu(0)=1$, $\mu(10)=0$, and
$\nu(0)=1-\varepsilon$, $\nu(10)=\varepsilon$.
The row of a coupling indexed by $10$ must be zero. Its column sums
then force the unique coupling to be
\[
 \pi(0,0)=1-\varepsilon,\qquad
 \pi(0,10)=\varepsilon,\qquad
 \pi(10,0)=\pi(10,10)=0.
\]
Its mean distance is $D=10\varepsilon$.

For the 1-Lipschitz function $f(x)=x$, center $a=0$, and $t=0$, the
original strict tail has mass zero. The target strict tail is
\[
 \nu\bigl(\{y:|f(y)|>r\}\bigr)
 =\begin{cases}
   \varepsilon,&0\le r<10,\\
   0,&r\ge10.
  \end{cases}
\]
For $0\le r<10$, the exceptional coupling mass is exactly
$\varepsilon$, so Theorem~\ref{thm:transfer} with
$\eta=\varepsilon$ is attained. In particular, a positive mean distance
does not imply exact containment in an interval $[-r,r]$ with $r<10$.
Even as $\varepsilon$ tends to zero and $D$ becomes arbitrarily small,
there remains positive mass at the distant point for each
$\varepsilon>0$. At $r=10$ the tail vanishes: the point with value $10$
lies on the boundary and is excluded by the strict inequality.

The same coupling proves the tail estimate separately for every
1-Lipschitz observable. It carries the chosen center through the
comparison, but supplies no guarantee that this center is a median
under $\nu$. All pairs and all distances in the argument belong to one
common metric space; comparison of different metric spaces would
require additional data specifying how their points are related.

\end{document}
